% =====================================================================
% Tutorial: Generador de Horarios a partir de un AST en Common Lisp
% =====================================================================
% Compilar con: pdflatex tutorial.tex
\documentclass[12pt,a4paper]{article}

\usepackage[utf8]{inputenc}
\usepackage[T1]{fontenc}
\usepackage[spanish]{babel}
\usepackage{amsmath,amssymb}
\usepackage{listings}
\usepackage{xcolor}
\usepackage{geometry}
\usepackage{hyperref}
\usepackage{enumitem}
\usepackage{booktabs}
\usepackage{caption}
\usepackage{fancyhdr}
\usepackage{titlesec}
\usepackage{parskip}

\geometry{left=3cm, right=2.5cm, top=3cm, bottom=3cm}
\setlength{\headheight}{14pt}

% -- Colores para listings ------------------------------------------
\definecolor{lisp-kw}{RGB}{0,0,160}
\definecolor{lisp-cmt}{RGB}{100,100,100}
\definecolor{lisp-str}{RGB}{160,0,0}
\definecolor{py-kw}{RGB}{0,0,180}
\definecolor{py-cmt}{RGB}{100,100,100}
\definecolor{py-str}{RGB}{160,0,0}
\definecolor{codebg}{RGB}{248,248,248}
\definecolor{frame-color}{RGB}{200,200,200}

% -- Estilo Common Lisp ---------------------------------------------
\lstdefinelanguage{CommonLisp}{
  morekeywords={defun,defvar,defmacro,defclass,let,let*,when,unless,
                dolist,loop,if,cond,case,typecase,format,progn,
                load,eval-when,setf,pushnew,mapcar,append,list,
                defclass*,def-entidad,def-comb,def-consulta,
                def-nppq,def-tqpq,def-tdpq,def-sbqp,
                def-minimizar,def-maximizar,
                op-and,op-or,op-not,
                op-igual,op-distinto,
                op-mayor,op-menor,
                op-mayor-igual,op-menor-igual,
                generar-python},
  sensitive=false,
  morecomment=[l]{;;},
  morestring=[b]",
}

\lstdefinestyle{lisp}{
  language=CommonLisp,
  backgroundcolor=\color{codebg},
  basicstyle=\ttfamily\small,
  keywordstyle=\color{lisp-kw}\bfseries,
  commentstyle=\color{lisp-cmt}\itshape,
  stringstyle=\color{lisp-str},
  frame=single,
  rulecolor=\color{frame-color},
  breaklines=true,
  breakatwhitespace=true,
  tabsize=2,
  showstringspaces=false,
  numbers=left,
  numberstyle=\tiny\color{lisp-cmt},
  numbersep=8pt,
}

\lstdefinestyle{python}{
  language=Python,
  backgroundcolor=\color{codebg},
  basicstyle=\ttfamily\small,
  keywordstyle=\color{py-kw}\bfseries,
  commentstyle=\color{py-cmt}\itshape,
  stringstyle=\color{py-str},
  frame=single,
  rulecolor=\color{frame-color},
  breaklines=true,
  breakatwhitespace=true,
  tabsize=4,
  showstringspaces=false,
  numbers=left,
  numberstyle=\tiny\color{py-cmt},
  numbersep=8pt,
}

\lstdefinestyle{shell}{
  backgroundcolor=\color{black},
  basicstyle=\ttfamily\small\color{white},
  frame=single,
  rulecolor=\color{black},
  breaklines=true,
  numbers=none,
}

\lstdefinestyle{output}{
  backgroundcolor=\color{codebg},
  basicstyle=\ttfamily\footnotesize,
  frame=single,
  rulecolor=\color{frame-color},
  breaklines=true,
  numbers=none,
}

% -- Encabezado -----------------------------------------------------
\pagestyle{fancy}
\fancyhf{}
\fancyhead[L]{\small Tutorial: Generador de Horarios}
\fancyhead[R]{\small\thepage}
\renewcommand{\headrulewidth}{0.4pt}

\hypersetup{
  colorlinks=true,
  linkcolor=blue!70!black,
  urlcolor=blue!70!black,
  citecolor=blue!70!black,
}

% -- Caracteres acentuados en listings (pdfLaTeX) -------------------
\lstset{literate=
  {á}{{\'{a}}}1 {é}{{\'{e}}}1 {í}{{\'{i}}}1 {ó}{{\'{o}}}1 {ú}{{\'{u}}}1
  {Á}{{\'{A}}}1 {É}{{\'{E}}}1 {Í}{{\'{I}}}1 {Ó}{{\'{O}}}1 {Ú}{{\'{U}}}1
  {ñ}{{\~{n}}}1 {Ñ}{{\~{N}}}1 {ü}{{\"u}}1 {Ü}{{\"U}}1
}

% -- Macros de conveniencia -----------------------------------------
\newcommand{\code}[1]{\texttt{#1}}
\newcommand{\file}[1]{\texttt{\textit{#1}}}
\newcommand{\kw}[1]{\textbf{\texttt{#1}}}

% ======================================================================
\begin{document}

\begin{titlepage}
  \centering
  \vspace*{3cm}
  {\LARGE\bfseries Tutorial del Generador de Horarios\\[0.5em]
   a partir de un \textit{Abstract Syntax Tree}\\
   en Common Lisp\par}
  \vspace{1.5cm}
  {\large Guía completa para la definición de dominios,\\
   modelado matemático y ejecución de la metaheurística\par}
  \vspace{2cm}
  \rule{\linewidth}{0.5pt}\\[0.5em]
  {\normalsize
    Este documento describe el uso del generador de código Python
    desarrollado como parte de una tesis de grado.
    Cubre desde la instalación del entorno hasta la interpretación
    de los resultados, y puede emplearse para modelar cualquier tipo
    de problema de asignación de horarios.\par}
  \rule{\linewidth}{0.5pt}
  \vfill
  {\small \today}
\end{titlepage}

\tableofcontents
\newpage

% ======================================================================
\section{Introducción}
% ======================================================================

El presente tutorial describe el uso del \textbf{Generador de Horarios},
una herramienta que permite definir un problema de asignación de horarios
mediante un Árbol de Sintaxis Abstracta (AST) escrito en Common Lisp y
genera automáticamente código Python con:

\begin{enumerate}
  \item Las funciones de evaluación de cada restricción del problema.
  \item El modelo matemático formal del problema de optimización.
  \item Un algoritmo genético que construye y optimiza el horario.
\end{enumerate}

La filosofía del sistema es \textbf{separar la descripción del problema
de su resolución}: el usuario define qué entidades existen en su dominio,
qué condiciones deben cumplirse y qué datos concretos hay; el generador
produce el código ejecutable de forma automática.

El sistema ha sido diseñado para ser \textbf{genérico}: aunque el ejemplo
que acompaña este tutorial corresponde a la asignación de defensas de
tesis, la misma herramienta puede emplearse para horarios de clases,
turnos de personal, asignación de salas, calendarios deportivos u
cualquier otro problema de naturaleza similar.

% ======================================================================
\section{Requisitos previos}
\label{sec:requisitos}
% ======================================================================

\subsection{Software necesario}

Para utilizar el generador se necesitan los siguientes programas:

\begin{description}
  \item[Steel Bank Common Lisp (SBCL)] Intérprete de Common Lisp.
    Disponible en \url{https://www.sbcl.org/}.
    Versión mínima recomendada: 2.0.

  \item[Python 3.10 o superior] Para ejecutar el código generado.
    Disponible en \url{https://www.python.org/}.
    No se requieren paquetes externos; la implementación utiliza
    únicamente la biblioteca estándar de Python.
\end{description}

\subsection{Instalación de SBCL}

La instalación varía según el sistema operativo.

\subsubsection*{Windows}

\begin{enumerate}
  \item Descargar el instalador desde \url{https://www.sbcl.org/platform-table.html}
        (seleccionar \textit{Windows x64}).
  \item Ejecutar el instalador y seguir los pasos predeterminados.
        SBCL se instala típicamente en
        \code{C:\textbackslash Program Files\textbackslash Steel Bank Common Lisp\textbackslash}.
  \item Para ejecutar \code{sbcl} desde PowerShell sin indicar la ruta completa,
        añadir el directorio de instalación al \code{PATH} del sistema
        (ejecutar una sola vez y reiniciar la terminal):
\end{enumerate}

\begin{lstlisting}[style=shell]
# PowerShell
[Environment]::SetEnvironmentVariable(
    "PATH",
    $env:PATH + ";C:\Program Files\Steel Bank Common Lisp",
    "User"
)
\end{lstlisting}

\subsubsection*{macOS}

En macOS la forma más sencilla es utilizar \textit{Homebrew}:

\begin{lstlisting}[style=shell]
# Terminal (bash / zsh)
brew install sbcl
\end{lstlisting}

Alternativamente puede descargarse el binario precompilado desde
\url{https://www.sbcl.org/platform-table.html} (seleccionar \textit{macOS}).
Tras la descarga, descomprimir y ejecutar el script de instalación
incluido:

\begin{lstlisting}[style=shell]
tar -xzf sbcl-<version>-x86-64-darwin-binary.tar.bz2
cd sbcl-<version>-x86-64-darwin
sudo sh install.sh
\end{lstlisting}

\subsubsection*{Linux}

En distribuciones basadas en \textit{Debian/Ubuntu}:

\begin{lstlisting}[style=shell]
# bash
sudo apt-get update
sudo apt-get install sbcl
\end{lstlisting}

En distribuciones basadas en \textit{Fedora/RHEL/Rocky Linux}:

\begin{lstlisting}[style=shell]
sudo dnf install sbcl
\end{lstlisting}

En \textit{Arch Linux}:

\begin{lstlisting}[style=shell]
sudo pacman -S sbcl
\end{lstlisting}

En todos los casos, para verificar que la instalación fue exitosa:

\begin{lstlisting}[style=shell]
sbcl --version
\end{lstlisting}

\subsection{Instalación de Python}

Python 3.10 o superior está disponible en \url{https://www.python.org/downloads/}.
En macOS y Linux es frecuente que ya esté instalado; puede verificarse con:

\begin{lstlisting}[style=shell]
python3 --version
\end{lstlisting}

En Windows, Python puede instalarse desde \url{https://www.python.org/} o
desde la Microsoft Store. Durante la instalación se recomienda marcar la
opción \textit{Add Python to PATH} para poder invocarlo simplemente como
\code{python} desde cualquier terminal.

\subsection{Obtención del código}

El código fuente se encuentra en el repositorio:
\begin{center}
  \url{https://github.com/salmaf04/restrictions_language}
\end{center}
Solo es necesaria la carpeta \file{generador-horarios/} del repositorio.
Puede clonarse el repositorio completo y trabajar desde esa carpeta, o
descargarse únicamente dicha carpeta. La estructura esperada del proyecto
es la siguiente:

\begin{verbatim}
generador-horarios/
+-- lib/
|   +-- clases-macros.lisp       <- infraestructura del AST (no modificar)
|   +-- generador.lisp           <- punto de entrada del generador
|   +-- gen-utils.lisp
|   +-- gen-ast-python.lisp
|   +-- gen-ast-math.lisp
|   +-- gen-serializar.lisp
|   +-- gen-consultas.lisp
|   +-- gen-restricciones.lisp
|   +-- gen-visualizacion.lisp
|   \-- gen-metaheuristica.lisp
\-- mi-dominio/
    \-- dominio.lisp             <- archivo que el usuario crea
\end{verbatim}

Los archivos bajo \file{lib/} son la \textbf{biblioteca del generador} y
no deben modificarse. El usuario únicamente necesita crear su propio
archivo de dominio (véase la sección~\ref{sec:dominio}).

% ======================================================================
\section{Conceptos fundamentales}
% ======================================================================

\subsection{El Árbol de Sintaxis Abstracta (AST)}

El AST es la representación interna del problema. Cuando el usuario
escribe definiciones en Common Lisp usando las macros del sistema, estas
construyen automáticamente objetos en memoria que forman un árbol. El
generador recorre ese árbol y produce el código Python equivalente.

El usuario \textbf{nunca manipula el árbol directamente}: trabaja con las
macros de alto nivel descritas en este documento.

\subsection{Tipos de cuantificadores de restricción}

El sistema reconoce cuatro tipos de cuantificadores para expresar
restricciones, más dos operadores de optimización:

\begin{table}[ht]
\centering
\caption{Cuantificadores del sistema}
\begin{tabular}{@{}lll@{}}
\toprule
\textbf{Macro} & \textbf{Nombre} & \textbf{Semántica} \\
\midrule
\code{def-nppq} & No Puede Pasar Que & Restricción dura.
  La condición nunca puede ser verdadera. \\
\code{def-tqpq} & Todo Que Pasa Que & Restricción dura.
  Toda entidad que satisface la premisa debe satisfacer la conclusión. \\
\code{def-tdpq} & Tampoco Debe Pasar Que & Restricción débil (blanda).
  Equivalente semántico a \code{nppq} pero con menor peso. \\
\code{def-sbqp} & Sería Bueno Que Pasara & Restricción débil (blanda).
  Equivalente semántico a \code{tqpq} pero con menor peso. \\
\code{def-minimizar} & Minimizar & Objetivo de optimización.
  Minimiza el máximo valor de una consulta sobre un dominio. \\
\code{def-maximizar} & Maximizar & Objetivo de optimización.
  Maximiza el mínimo valor de una consulta sobre un dominio. \\
\bottomrule
\end{tabular}
\end{table}

Las restricciones \textbf{duras} (\code{nppq}, \code{tqpq}) reciben un
peso alto (por defecto 1000) en la función objetivo. Las
\textbf{blandas} (\code{tdpq}, \code{sbqp}, \code{minimizar},
\code{maximizar}) reciben un peso bajo (por defecto 1). La
metaheurística minimiza la suma ponderada de violaciones.

\subsection{Operadores lógicos disponibles}

Para construir condiciones se dispone de los siguientes operadores:

\begin{table}[ht]
\centering
\caption{Operadores lógicos y relacionales}
\begin{tabular}{@{}ll@{}}
\toprule
\textbf{Macro} & \textbf{Significado} \\
\midrule
\code{op-and e1 e2 \ldots} & Conjunción ($e_1 \wedge e_2 \wedge \cdots$) \\
\code{op-or  e1 e2 \ldots} & Disyunción ($e_1 \vee e_2 \vee \cdots$) \\
\code{op-not e}            & Negación ($\neg e$) \\
\code{op-igual a b}        & Igualdad ($a = b$) \\
\code{op-distinto a b}     & Desigualdad ($a \neq b$) \\
\code{op-mayor a b}        & Mayor estricto ($a > b$) \\
\code{op-menor a b}        & Menor estricto ($a < b$) \\
\code{op-mayor-igual a b}  & Mayor o igual ($a \geq b$) \\
\code{op-menor-igual a b}  & Menor o igual ($a \leq b$) \\
\bottomrule
\end{tabular}
\end{table}

% ======================================================================
\section{Definición del dominio}
\label{sec:dominio}
% ======================================================================

Todo el dominio se describe en un único archivo Lisp que el usuario crea.
Este archivo debe:

\begin{enumerate}
  \item Cargar la biblioteca del generador.
  \item Definir los \textit{mixins} (atributos compartidos).
  \item Definir las entidades del dominio.
  \item Definir la entidad combinación (la ``fila'' del horario).
  \item Definir las consultas (auxiliares y con iteración).
  \item Definir las restricciones.
  \item Definir los datos concretos del problema.
  \item Llamar a \code{generar-python}.
\end{enumerate}

La estructura general del archivo es la siguiente:

\begin{lstlisting}[style=lisp, caption={Estructura del archivo de dominio}]
;; 1. Cargar la biblioteca
(load (merge-pathnames "lib/clases-macros.lisp" (truename ".")))
(load (merge-pathnames "lib/generador.lisp"     (truename ".")))

;; 2. Mixins
(defclass* tiene-nombre () (nombre))

;; 3. Entidades
(def-entidad recurso-a () (atributo1 atributo2))

;; 4. Entidad combinación
(def-comb asignacion (identidad variable1 variable2))

;; 5. Consultas
(def-consulta mi-consulta (arg1 arg2) ...)

;; 6. Restricciones
(defvar restriccion-1 (def-nppq ...))

;; 7. Datos concretos
(def-recurso-a inst-1 valor1 valor2)

;; 8. Generar
(generar-python "salida/horario.py"
  :entidad-horario 'asignacion
  ...)

(sb-ext:exit)
\end{lstlisting}

Las subsecciones siguientes describen cada bloque en detalle.

\subsection{Carga de la biblioteca}

El archivo debe comenzar cargando los dos archivos de la biblioteca.
La ruta es relativa al directorio desde el que se ejecutará SBCL,
que debe ser la raíz del proyecto (\file{generador-horarios/}):

\begin{lstlisting}[style=lisp, caption={Carga de la biblioteca}]
(load (merge-pathnames "lib/clases-macros.lisp" (truename ".")))
(load (merge-pathnames "lib/generador.lisp"     (truename ".")))
\end{lstlisting}

\subsection{Definición de \textit{mixins}}

Un \textit{mixin} es un conjunto de atributos que pueden compartir varias
entidades. Se define con la macro \code{defclass*}:

\begin{lstlisting}[style=lisp, caption={Sintaxis de defclass*}]
(defclass* nombre-mixin (padres...) (slot1 slot2 ...))
\end{lstlisting}

El ejemplo más común es un mixin con un atributo \code{nombre}:

\begin{lstlisting}[style=lisp, caption={Ejemplo de mixin}]
(defclass* tiene-nombre () (nombre))
\end{lstlisting}

Tras esta definición, cualquier entidad que herede de \code{tiene-nombre}
dispondrá automáticamente del atributo \code{nombre} y de la función de
acceso \code{(nombre entidad)}.

\subsection{Definición de entidades}
\label{sec:entidades}

Una entidad representa un tipo de objeto del dominio. Se define con
\code{def-entidad}:

\begin{lstlisting}[style=lisp, caption={Sintaxis de def-entidad}]
(def-entidad nombre-entidad (padres-mixin...) (slot1 slot2 ...))
\end{lstlisting}

\begin{description}
  \item[\code{nombre-entidad}] Nombre del tipo de entidad (símbolo Lisp).
  \item[\code{(padres-mixin...)}] Lista de \textit{mixins} de los que
    hereda. Puede estar vacía: \code{()}.
  \item[\code{(slot1 slot2 ...)}] Atributos propios de esta entidad.
    También puede estar vacía si la entidad solo tiene atributos heredados.
\end{description}

\textbf{Ejemplo}: definir un tipo \code{profesor} que tiene nombre
(heredado del mixin) y grado académico (propio):

\begin{lstlisting}[style=lisp, caption={Definición de la entidad profesor}]
(def-entidad profesor (tiene-nombre) (grado))
\end{lstlisting}

Internamente, \code{def-entidad} genera:
\begin{itemize}
  \item Una clase CLOS para las instancias concretas.
  \item Una clase CLOS para la referencia genérica (usada en restricciones).
  \item Una macro \code{def-profesor} para crear instancias.
  \item Funciones de acceso AST: \code{(nombre prof)}, \code{(grado prof)}.
\end{itemize}

\subsubsection{Entidades enumeradas}

Una \textbf{entidad enumerada} es aquella cuyas instancias son fijas y
conocidas de antemano (como grados académicos, días de la semana o
estados). Se define igual que cualquier entidad, pero generalmente con
un único atributo identificador:

\begin{lstlisting}[style=lisp, caption={Entidad enumerada: grado académico}]
(def-entidad grado () (id))

(def-grado grado-lic "Lic")
(def-grado grado-msc "Msc")
(def-grado grado-dr  "Dr")
\end{lstlisting}

\subsection{Definición de la entidad combinación}
\label{sec:comb}

La entidad combinación representa una \textbf{fila del horario}: qué
identidad (p.ej., un tribunal, un grupo, un empleado) es asignada a qué
recursos (p.ej., fecha, hora, aula). Se define con \code{def-comb}:

\begin{lstlisting}[style=lisp, caption={Sintaxis de def-comb}]
(def-comb nombre-combinacion (slot-identidad variable1 variable2 ...))
\end{lstlisting}

\begin{description}
  \item[\code{slot-identidad}] El atributo que identifica de forma única
    \textit{qué} se está asignando. Hay exactamente una instancia de la
    combinación por cada valor de este atributo.
  \item[\code{variable1, variable2, ...}] Los atributos que el algoritmo
    genético decide: cuándo, dónde, en qué condiciones se realiza la
    asignación.
\end{description}

\textbf{Ejemplo} para un horario de defensas de tesis:

\begin{lstlisting}[style=lisp, caption={Entidad combinación para defensas de tesis}]
(def-comb asignacion (tribunal fecha momento local))
\end{lstlisting}

Aquí, \code{tribunal} es la identidad (hay una defensa por tribunal) y
\code{fecha}, \code{momento}, \code{local} son las variables de decisión
que el algoritmo optimizará.

\subsection{Definición de consultas}
\label{sec:consultas}

Las consultas son expresiones con nombre y parámetros que pueden
reutilizarse dentro de restricciones o de otras consultas.
Se definen con \code{def-consulta} en dos variantes: con o sin
iteración sobre las instancias del horario.

\subsubsection{Consulta sin iteración}

Una consulta sin iteración recibe argumentos explícitos y devuelve
una expresión del AST construida con los operadores disponibles.
Es útil para encapsular condiciones o sub-expresiones que se repiten
en varias restricciones, aunque su uso es opcional: una restricción
puede construir directamente su expresión sin necesidad de una consulta
auxiliar.

\begin{lstlisting}[style=lisp, caption={Sintaxis de consulta sin iteración}]
(def-consulta nombre-consulta (arg1 arg2 ...)
  :comprueba
  <expresion-del-AST>)
\end{lstlisting}

\textbf{Ejemplo}: expresión que verifica si un profesor pertenece al
tribunal de una asignación dada:

\begin{lstlisting}[style=lisp, caption={Consulta auxiliar: profesor en tribunal}]
(def-consulta profesor-en-tribunal (prof asig)
  :comprueba
  (op-or
    (op-igual (tutor      (tribunal asig)) prof)
    (op-igual (oponente   (tribunal asig)) prof)
    (op-igual (presidente (tribunal asig)) prof)
    (op-igual (vocal      (tribunal asig)) prof)
    (op-igual (secretario (tribunal asig)) prof)))
\end{lstlisting}

El acceso encadenado \code{(tutor (tribunal asig))} navega los atributos
del AST: primero obtiene el \code{tribunal} de la asignación \code{asig}
y luego su atributo \code{tutor}. El generador produce el código de
acceso equivalente en Python de forma automática.

Esta consulta puede usarse como subexpresión dentro de otras consultas
o restricciones.

\subsubsection{Consulta con iteración}

Una consulta con iteración recorre uno o varios dominios de entidades,
filtra con una condición y acumula un resultado numérico. Es el
mecanismo principal para calcular métricas sobre el horario completo:
conteos, sumas, o cualquier agregado expresable con las operaciones
disponibles.

\begin{lstlisting}[style=lisp, caption={Sintaxis de consulta con iteración}]
(def-consulta nombre-consulta (args-extra...)
  :itera-sobre (var1 entidad1) (var2 entidad2) ...
  :comprueba   <condicion-booleana>
  :operacion   <tipo-de-acumulacion>
  :devuelve    nombre-resultado)
\end{lstlisting}

\begin{description}
  \item[\code{:itera-sobre}] Define los bucles anidados. Cada par
    \code{(variable entidad)} recorre todas las instancias de ese tipo
    definidas en el dominio.
  \item[\code{:comprueba}] Condición booleana que determina qué
    combinaciones de variables se acumulan en el resultado.
  \item[\code{:operacion}] Tipo de acumulación:
    \begin{itemize}
      \item \code{suma} o \code{contar}: cuenta cuántas veces se cumple
            la condición.
      \item \code{contar-distintos-dias}: cuenta cuántos valores distintos
            del atributo \code{fecha} aparecen entre los elementos que
            cumplen la condición.
    \end{itemize}
  \item[\code{:devuelve}] Nombre descriptivo del resultado (informativo,
    no afecta al código generado).
\end{description}

\textbf{Ejemplo}: contar cuántas combinaciones de profesor y par de
asignaciones presentan conflicto de horario:

\begin{lstlisting}[style=lisp, caption={Consulta con iteración: conflictos de horario}]
(def-consulta conflictos-de-horario ()
  :itera-sobre (p profesor) (a1 asignacion) (a2 asignacion)
  :comprueba
  (op-and
    (profesor-en-tribunal p a1)
    (profesor-en-tribunal p a2)
    (op-igual    (fecha    a1) (fecha    a2))
    (op-igual    (momento  a1) (momento  a2))
    (op-distinto (tribunal a1) (tribunal a2)))
  :operacion suma
  :devuelve  n-conflictos)
\end{lstlisting}

\textbf{Ejemplo}: contar cuántos días distintos asiste un profesor
concreto a lo largo del horario:

\begin{lstlisting}[style=lisp, caption={Consulta: días distintos de asistencia}]
(def-consulta dias-asistidos (prof)
  :itera-sobre (asig asignacion)
  :comprueba   (profesor-en-tribunal prof asig)
  :operacion   contar-distintos-dias
  :devuelve    n-dias)
\end{lstlisting}

Las consultas con iteración son las que deben listarse en el parámetro
\code{:consultas} de \code{generar-python}. Las consultas sin iteración
se referencian automáticamente desde las que las utilizan y no necesitan
declararse explícitamente en esa lista.

\subsection{Definición de restricciones}
\label{sec:restricciones}

Las restricciones expresan condiciones sobre el horario y se almacenan
en variables globales con \code{defvar}. Cada restricción envuelve una
expresión del AST, que puede ser una llamada a una consulta, una
combinación de operadores lógicos o cualquier otra construcción válida.
El generador traduce cada restricción a una función Python que devuelve
un valor numérico: cuanto mayor sea ese valor, mayor es el
incumplimiento.

\subsubsection{Restricción nppq (No Puede Pasar Que)}

Expresa que la expresión dada \textbf{no debe ocurrir o debe valer
cero}. El generador produce una función cuyo resultado se multiplica
por el peso de la restricción en la función objetivo; un valor de cero
indica cumplimiento total.

\begin{lstlisting}[style=lisp, caption={Sintaxis de restricción nppq}]
(defvar nombre-restriccion
  (def-nppq <expresion-del-AST>))
\end{lstlisting}

\textbf{Ejemplo}: prohibir cualquier conflicto de horario entre
profesores:

\begin{lstlisting}[style=lisp, caption={Restricción nppq: sin conflictos de horario}]
(defvar restriccion-sin-conflictos-de-horario
  (def-nppq (op-mayor conflictos-de-horario 0)))
\end{lstlisting}

\subsubsection{Restricción tqpq (Todo Que Pasa Que)}

Expresa que la expresión dada \textbf{debe cumplirse}. La función
generada devuelve un valor positivo proporcional al grado de
incumplimiento, y cero cuando la condición se satisface plenamente.

\begin{lstlisting}[style=lisp, caption={Sintaxis de restricción tqpq}]
(defvar nombre-restriccion
  (def-tqpq <expresion-del-AST>))
\end{lstlisting}

\subsubsection{Restricción débil: minimizar}

Expresa un objetivo de optimización: minimizar el valor máximo de una
consulta evaluada sobre cada instancia de un dominio de entidades.
Es útil para equilibrar cargas o recursos entre los distintos elementos
del problema.

\begin{lstlisting}[style=lisp, caption={Restricción débil: minimizar carga máxima}]
(defvar restriccion-minimizar-carga-maxima
  (def-minimizar dias-asistidos profesor))
\end{lstlisting}

El primer argumento es la consulta a evaluar y el segundo es el tipo
de entidad sobre el que se itera para calcular el máximo.

\subsubsection{Restricciones blandas: tdpq y sbqp}

\code{def-tdpq} (Tampoco Debe Pasar Que) es la versión blanda de
\code{nppq}: expresa algo que preferiblemente no debe ocurrir, pero
cuya violación tiene un coste menor en la función objetivo.

\code{def-sbqp} (Sería Bueno Que Pasara) es la versión blanda de
\code{tqpq}: expresa algo que sería deseable que se cumpliera, sin
que su incumplimiento sea crítico.

Ambas tienen la misma sintaxis que sus equivalentes duros y aceptan
el mismo tipo de expresiones del AST:

\begin{lstlisting}[style=lisp, caption={Restricciones blandas}]
(defvar restriccion-blanda-1
  (def-tdpq <expresion-que-se-prefiere-no-ocurra>))

(defvar restriccion-blanda-2
  (def-sbqp <expresion-deseable>))
\end{lstlisting}

\subsection{Definición de datos concretos}
\label{sec:datos}

Una vez definidas las entidades, se crean las instancias concretas del
problema usando las macros generadas automáticamente por \code{def-entidad}.
La macro tiene la forma \code{def-<nombre-entidad>}:

\begin{lstlisting}[style=lisp, caption={Sintaxis general para instanciar entidades}]
(def-<nombre-entidad> nombre-instancia valor1 valor2 ...)
\end{lstlisting}

Los valores se asignan en orden: primero los atributos heredados de los
\textit{mixins} (en el orden en que se declararon los mixins), y luego
los atributos propios.

\textbf{Ejemplo completo de datos para el horario de tesis}:

\begin{lstlisting}[style=lisp, caption={Datos concretos del dominio de tesis}]
;; Grados (enumeración)
(def-grado grado-lic "Lic")
(def-grado grado-msc "Msc")
(def-grado grado-dr  "Dr")

;; Profesores: (nombre grado)
;; Orden: slot heredado 'nombre' primero, luego slot propio 'grado'
(def-profesor prof-piad   "Alejandro Piad"  grado-dr)
(def-profesor prof-suarez "Carlos Suarez"   grado-msc)
(def-profesor prof-garcia "Maria Garcia"    grado-dr)

;; Estudiantes: (nombre)
(def-estudiante est-rodriguez "Laura Rodriguez")
(def-estudiante est-fernandez "Carlos Fernandez")

;; Tribunales: (estudiante tutor oponente presidente vocal secretario)
(def-tribunal trib-1
  est-rodriguez prof-piad prof-suarez prof-garcia prof-piad prof-suarez)

;; Fechas disponibles: (dia mes)
(def-fecha fecha-1 1 6)
(def-fecha fecha-2 2 6)
(def-fecha fecha-3 3 6)

;; Momentos disponibles: (hora)
(def-momento momento-manana   9)
(def-momento momento-mediodia 11)
(def-momento momento-tarde    14)

;; Locales disponibles: (nombre)
(def-local local-a "Sala A")
(def-local local-b "Sala B")
\end{lstlisting}

\textbf{Importante}: el orden de los argumentos en \code{def-<entidad>}
debe coincidir exactamente con el orden en que se declararon los slots
en la definición de la entidad (primero los heredados, luego los propios).

Todas las instancias se registran automáticamente en un repositorio
interno. El generador las lee para construir las constantes del módulo
Python y poblar los pools del algoritmo genético.

% ======================================================================
\section{Llamada al generador}
\label{sec:llamada}
% ======================================================================

Una vez definido el dominio completo, se llama a la función
\code{generar-python} para producir el archivo Python:

\begin{lstlisting}[style=lisp, caption={Llamada completa a generar-python}]
(generar-python "mi-dominio/horario.py"
  ;; Entidad combinación (la fila del horario)
  :entidad-horario   'asignacion

  ;; Slot que identifica cada fila (hay exactamente una por valor)
  :slot-identidad    'tribunal

  ;; Slots que el algoritmo genético optimizará
  :slots-variables   '(fecha momento local)

  ;; Consultas con :itera-sobre (defvar, no defun)
  :consultas         '(conflictos-de-horario dias-asistidos)

  ;; Restricciones duras (nppq / tqpq)
  :restricciones-duras   '(restriccion-sin-conflictos-de-horario)

  ;; Restricciones blandas (tdpq / sbqp / minimizar / maximizar)
  :restricciones-blandas '(restriccion-minimizar-carga-maxima)

  ;; Pesos del modelo matemático (opcionales, defaults: 1000 y 1)
  :peso-dura   1000
  :peso-blanda    1)
\end{lstlisting}

\subsection{Descripción de cada parámetro}

\begin{description}
  \item[\code{:entidad-horario}] Símbolo de la entidad combinación definida
    con \code{def-comb}. El horario es una lista de instancias de este tipo.

  \item[\code{:slot-identidad}] Nombre del atributo de la combinación que
    actúa como clave: hay exactamente una asignación por cada valor distinto
    de este slot. El algoritmo genético \textit{no modifica} este slot.

  \item[\code{:slots-variables}] Lista de atributos que el algoritmo
    genético puede cambiar. Cada uno se elige aleatoriamente de entre las
    instancias del tipo correspondiente que hayan sido definidas con
    \code{def-<entidad>}.

  \item[\code{:consultas}] Lista de las variables globales que almacenan
    nodos \code{consulta-simple} con \code{:itera-sobre}. Estas son las
    que se convierten en funciones Python. Las consultas auxiliares (sin
    \code{:itera-sobre}, definidas como \code{defun}) se referencian
    automáticamente desde las que sí iteran.

  \item[\code{:restricciones-duras}] Lista de variables globales con
    cuantificadores \code{nppq} o \code{tqpq}. Su violación tiene un
    coste alto en la función objetivo (\code{:peso-dura}).

  \item[\code{:restricciones-blandas}] Lista de variables globales con
    cuantificadores \code{tdpq}, \code{sbqp}, \code{minimizar} o
    \code{maximizar}. Su violación tiene un coste bajo (\code{:peso-blanda}).

  \item[\code{:peso-dura}] Peso numérico para las restricciones duras.
    Por defecto: 1000.

  \item[\code{:peso-blanda}] Peso numérico para las restricciones blandas.
    Por defecto: 1.
\end{description}

\subsection{Finalización del script}

Al final del archivo de dominio, se añade \code{(sb-ext:exit)} para que
SBCL termine limpiamente tras generar el archivo:

\begin{lstlisting}[style=lisp]
(sb-ext:exit)
\end{lstlisting}

% ======================================================================
\section{Ejecución}
\label{sec:ejecucion}
% ======================================================================

\subsection{Paso 1: generar el archivo Python}

Situarse en la raíz del proyecto (\file{generador-horarios/}) y ejecutar
SBCL con el archivo de dominio. Los comandos varían según el sistema
operativo.

\subsubsection*{Windows (PowerShell)}

Si SBCL fue añadido al \code{PATH} (véase sección~\ref{sec:requisitos}):

\begin{lstlisting}[style=shell]
cd "C:\ruta\al\generador-horarios"
sbcl --script mi-dominio/dominio.lisp
\end{lstlisting}

Si no fue añadido al \code{PATH}, usar la ruta completa:

\begin{lstlisting}[style=shell]
cd "C:\ruta\al\generador-horarios"
& "C:\Program Files\Steel Bank Common Lisp\sbcl.exe" --script mi-dominio/dominio.lisp
\end{lstlisting}

\subsubsection*{macOS y Linux (Terminal)}

\begin{lstlisting}[style=shell]
cd /ruta/al/generador-horarios
sbcl --script mi-dominio/dominio.lisp
\end{lstlisting}

En todos los casos, un mensaje de confirmación indicará la ruta del
archivo generado:

\begin{lstlisting}[style=output]
Python generado en: mi-dominio/horario.py
\end{lstlisting}

\subsection{Paso 2: ejecutar el algoritmo genético}

\subsubsection*{Windows (PowerShell)}

\begin{lstlisting}[style=shell]
python mi-dominio/horario.py
\end{lstlisting}

Si \code{python} no está en el \code{PATH}:

\begin{lstlisting}[style=shell]
& "C:\Users\<usuario>\AppData\Local\Programs\Python\Python313\python.exe" `
    mi-dominio/horario.py
\end{lstlisting}

\subsubsection*{macOS y Linux (Terminal)}

\begin{lstlisting}[style=shell]
python3 mi-dominio/horario.py
\end{lstlisting}

% ======================================================================
\section{Interpretación de la salida}
\label{sec:salida}
% ======================================================================

\subsection{Modelo de optimización}

La primera sección impresa corresponde al modelo matemático abstracto:

\begin{lstlisting}[style=output]
============================================================
MODELO DE OPTIMIZACION
============================================================

INDICES Y CONJUNTOS:
  T  =  conjunto de tribunals      (|T| = 3)
  F  =  conjunto de fecha disponibles   (|F| = 3)
  M  =  conjunto de momento disponibles (|M| = 3)
  L  =  conjunto de local disponibles   (|L| = 2)
  P  =  conjunto de profesor            (|P| = 6)

VARIABLES DE DECISION:
  Para cada tribunal t en T se asigna:
    fecha_t   in F   (fecha)
    momento_t in M   (momento)
    local_t   in L   (local)

FUNCION OBJETIVO:
  minimizar z(h) donde:
    z(h) = 1000 * sum_i incumpl_i(h)   [duras]
         +    1 * sum_j incumpl_j(h)   [blandas]

RESTRICCIONES DURAS  (nppq / tqpq):
  R1  [SIMPLE / DURA  (nppq) -- No Puede Pasar Que]:
    conflictos_de_horario(h) = 0
    donde:
      conflictos_de_horario(h) = Sum_{p in P, a1 in A, a2 in A}
        [prof_en_tribunal(p,a1) /\ prof_en_tribunal(p,a2)
         /\ fecha(a1)=fecha(a2) /\ momento(a1)=momento(a2)
         /\ tribunal(a1) != tribunal(a2)]

RESTRICCIONES BLANDAS  (minimizar):
  R2  [Minimizar]:
    minimizar max_{x in P}  dias_asistidos(h, x)
\end{lstlisting}

\subsubsection{Lectura del modelo}

\begin{description}
  \item[Índices y conjuntos] Muestra los tipos de entidades involucrados
    y cuántas instancias concretas hay de cada uno en el problema.

  \item[Variables de decisión] Indica qué atributos se están decidiendo
    para cada valor del slot identidad. La notación $v_t$ significa
    ``el valor del atributo $v$ asignado al tribunal $t$''.

  \item[Función objetivo] Fórmula de minimización con los pesos configurados.

  \item[Restricciones] Para cada restricción se muestra su tipo,
    la expresión matemática que debe satisfacerse y, cuando aplica,
    la definición formal de las consultas que intervienen.
\end{description}

\subsection{Progreso del algoritmo genético}

Durante la ejecución se imprime el progreso cada 50 generaciones:

\begin{lstlisting}[style=output]
Generacion    0: costo inicial = 3000
Generacion   50: mejor costo   = 1000
Generacion  100: mejor costo   = 1
Generacion  150: mejor costo   = 1
...
\end{lstlisting}

Un costo de 0 indica que todas las restricciones duras se cumplen
perfectamente y los objetivos de optimización han alcanzado su mínimo.
Si el algoritmo encuentra una solución perfecta antes de agotar las
generaciones, lo indica explícitamente.

\subsection{Horario propuesto y evaluación}

Al finalizar, se imprime el horario concreto elegido por la metaheurística
y la evaluación de cada restricción:

\begin{lstlisting}[style=output]
============================================================
HORARIO PROPUESTO POR LA METAHEURISTICA
============================================================

  Defensa 1:
    TRIBUNAL : (estudiante=Laura Rodriguez, tutor=Alejandro Piad, ...)
    FECHA    : (dia=2, mes=6)
    MOMENTO  : 11
    LOCAL    : Sala A

  Defensa 2:
    TRIBUNAL : (estudiante=Carlos Fernandez, tutor=Maria Garcia, ...)
    FECHA    : (dia=2, mes=6)
    MOMENTO  : 14
    LOCAL    : Sala B

============================================================
EVALUACION DE RESTRICCIONES
============================================================
Restricciones simples / duras  (peso 1000):
  restriccion_sin_conflictos_de_horario: 0 incumplimientos  ->  costo: 0
Restricciones debiles / blandas  (peso 1):
  restriccion_minimizar_carga_maxima: 1  ->  costo: 1
Costo total: 1
============================================================
\end{lstlisting}

\textbf{Interpretación}:
\begin{itemize}
  \item \textit{0 incumplimientos} en la restricción dura significa que
    ningún profesor debe asistir a dos defensas simultáneas.
  \item El valor 1 en la restricción blanda indica que algún profesor
    asiste al menos un día; dado el tamaño del problema esto es inevitable.
  \item El costo total de 1 (solo blandas) es la mejor solución posible
    para este conjunto de datos.
\end{itemize}

% ======================================================================
\section{Ejemplo completo: horario de defensas de tesis}
\label{sec:ejemplo}
% ======================================================================

Esta sección presenta el archivo de dominio completo utilizado como caso
de prueba del sistema.

\begin{lstlisting}[style=lisp, caption={Archivo completo: tesis/dominio.lisp}]
;; =====================================================================
;; DOMINIO: HORARIO DE DEFENSAS DE TESIS
;; =====================================================================
;; Ejecutar desde la carpeta generador-horarios/:
;;   sbcl --script tesis/dominio.lisp
;; =====================================================================

;; Cargar infraestructura
(load (merge-pathnames "lib/clases-macros.lisp" (truename ".")))
(load (merge-pathnames "lib/generador.lisp"     (truename ".")))


;; -- Mixins --------------------------------------------------------
(defclass* tiene-nombre () (nombre))


;; -- Entidades -----------------------------------------------------
(def-entidad grado     ()           (id))
(def-entidad profesor  (tiene-nombre) (grado))
(def-entidad estudiante (tiene-nombre))
(def-entidad tribunal  ()
  (estudiante tutor oponente presidente vocal secretario))
(def-entidad fecha     () (dia mes))
(def-entidad momento   () (hora))
(def-entidad local     (tiene-nombre))


;; -- Entidad combinación -------------------------------------------
(def-comb asignacion (tribunal fecha momento local))


;; -- Consultas -----------------------------------------------------

;; Auxiliar: profesor pertenece al tribunal de la asignacion
(def-consulta profesor-en-tribunal (prof asig)
  :comprueba
  (op-or
    (op-igual (tutor      (tribunal asig)) prof)
    (op-igual (oponente   (tribunal asig)) prof)
    (op-igual (presidente (tribunal asig)) prof)
    (op-igual (vocal      (tribunal asig)) prof)
    (op-igual (secretario (tribunal asig)) prof)))

;; Cuenta pares de defensas con conflicto de horario
(def-consulta conflictos-de-horario ()
  :itera-sobre (p profesor) (a1 asignacion) (a2 asignacion)
  :comprueba
  (op-and
    (profesor-en-tribunal p a1)
    (profesor-en-tribunal p a2)
    (op-igual    (fecha    a1) (fecha    a2))
    (op-igual    (momento  a1) (momento  a2))
    (op-distinto (tribunal a1) (tribunal a2)))
  :operacion suma
  :devuelve  n-conflictos)

;; Cuenta dias distintos que asiste un profesor
(def-consulta dias-asistidos (prof)
  :itera-sobre (asig asignacion)
  :comprueba   (profesor-en-tribunal prof asig)
  :operacion   contar-distintos-dias
  :devuelve    n-dias)


;; -- Restricciones -------------------------------------------------

;; DURA: ningun profesor en dos defensas al mismo tiempo
(defvar restriccion-sin-conflictos-de-horario
  (def-nppq (op-mayor conflictos-de-horario 0)))

;; BLANDA: minimizar la carga del profesor que mas dias asiste
(defvar restriccion-minimizar-carga-maxima
  (def-minimizar dias-asistidos profesor))


;; -- Datos concretos -----------------------------------------------

(def-grado grado-lic "Lic")
(def-grado grado-msc "Msc")
(def-grado grado-dr  "Dr")

(def-profesor prof-piad    "Alejandro Piad"  grado-dr)
(def-profesor prof-suarez  "Carlos Suarez"   grado-msc)
(def-profesor prof-garcia  "Maria Garcia"    grado-dr)
(def-profesor prof-torres  "Luis Torres"     grado-msc)
(def-profesor prof-mendez  "Ana Mendez"      grado-dr)
(def-profesor prof-herrera "Pedro Herrera"   grado-msc)

(def-estudiante est-rodriguez "Laura Rodriguez")
(def-estudiante est-fernandez "Carlos Fernandez")
(def-estudiante est-lopez     "Sofia Lopez")

(def-tribunal trib-1
  est-rodriguez prof-piad   prof-suarez  prof-garcia  prof-torres  prof-mendez)
(def-tribunal trib-2
  est-fernandez prof-garcia prof-herrera prof-piad    prof-suarez  prof-torres)
(def-tribunal trib-3
  est-lopez     prof-mendez prof-torres  prof-suarez  prof-herrera prof-garcia)

(def-fecha fecha-1 1 6)
(def-fecha fecha-2 2 6)
(def-fecha fecha-3 3 6)

(def-momento momento-manana   9)
(def-momento momento-mediodia 11)
(def-momento momento-tarde    14)

(def-local local-a "Sala A")
(def-local local-b "Sala B")


;; -- Generar -------------------------------------------------------
(generar-python "tesis/horario.py"
  :entidad-horario   'asignacion
  :slot-identidad    'tribunal
  :slots-variables   '(fecha momento local)
  :consultas         '(conflictos-de-horario dias-asistidos)
  :restricciones-duras   '(restriccion-sin-conflictos-de-horario)
  :restricciones-blandas '(restriccion-minimizar-carga-maxima)
  :peso-dura   1000
  :peso-blanda    1)

(sb-ext:exit)
\end{lstlisting}

% ======================================================================
\section{Generalización a otros tipos de horarios}
\label{sec:generalizacion}
% ======================================================================

El sistema no está ligado al dominio de defensas de tesis. Puede usarse
para cualquier problema donde se deba asignar recursos a identidades
minimizando conflictos. La siguiente tabla muestra algunos ejemplos:

\begin{table}[ht]
\centering
\caption{Ejemplos de aplicación del generador}
\begin{tabular}{@{}p{3cm}p{3.5cm}p{4cm}p{3cm}@{}}
\toprule
\textbf{Problema} & \textbf{Identidad} & \textbf{Variables} & \textbf{Restricciones típicas} \\
\midrule
Horario de clases & grupo de alumnos & día, hora, aula, profesor
  & Sin solapamiento de profesores; aula con capacidad suficiente \\
Turnos de personal & empleado & día, turno, puesto
  & Horas máximas semanales; descanso mínimo entre turnos \\
Calendarios deportivos & partido & jornada, hora, estadio
  & Equipos sin dos partidos en un mismo día; uso alternado del estadio \\
Asignación de aulas & evento & fecha, hora, sala
  & Sin colisión de salas; compatibilidad de capacidad \\
\bottomrule
\end{tabular}
\end{table}

Para adaptar el generador a un nuevo dominio, el proceso es siempre el mismo:

\begin{enumerate}
  \item \textbf{Identificar la entidad combinación}: qué es lo que se está
    asignando (el elemento principal del horario) y qué atributos tiene.
  \item \textbf{Separar identidad de variables}: determinar cuál atributo
    identifica unívocamente cada elemento (identidad) y cuáles son los
    que se asignan (variables).
  \item \textbf{Definir las restricciones}: expresar en términos de
    consultas cuándo se produce un conflicto o una situación indeseable.
  \item \textbf{Categorizar las restricciones}: decidir cuáles son duras
    (nunca pueden violarse) y cuáles son blandas (se prefiere no violarlas
    pero no es crítico).
  \item \textbf{Proporcionar los datos}: definir las instancias concretas
    de cada tipo de entidad.
  \item \textbf{Ejecutar el generador} y analizar los resultados.
\end{enumerate}

\subsection{Lista de verificación antes de ejecutar}

Antes de llamar a \code{generar-python}, conviene verificar lo siguiente:

\begin{itemize}[label=$\square$]
  \item La entidad combinación está definida con \code{def-comb} y
    contiene exactamente un slot de identidad más los slots variables.
  \item Cada slot variable tiene al menos dos instancias concretas
    definidas (de lo contrario el algoritmo no puede variar).
  \item Cada consulta pasada en \code{:consultas} fue definida con
    \code{:itera-sobre} (es un \code{defvar}, no un \code{defun}).
  \item Las restricciones pasadas en \code{:restricciones-duras} y
    \code{:restricciones-blandas} son variables globales (definidas
    con \code{defvar}).
  \item El directorio destino del archivo Python existe.
  \item SBCL se ejecuta desde la raíz del proyecto
    (\file{generador-horarios/}).
\end{itemize}

% ======================================================================
\section{Referencia rápida de macros}
\label{sec:referencia}
% ======================================================================

\begin{table}[ht]
\centering
\caption{Macros principales del sistema}
\begin{tabular}{@{}p{5cm}p{8cm}@{}}
\toprule
\textbf{Macro / Función} & \textbf{Uso} \\
\midrule
\code{defclass* N () (s1 s2)} & Define mixin con atributos \code{s1}, \code{s2} \\
\code{def-entidad N (mix) (s)} & Define entidad heredando mixins, con slot \code{s} \\
\code{def-comb N (id v1 v2)} & Define entidad combinación \\
\code{def-N inst v1 v2} & Crea instancia de entidad N \\
\code{def-consulta N (a) :comprueba e} & Consulta auxiliar (retorna expresión AST) \\
\code{def-consulta N () :itera-sobre ...} & Consulta con iteración (devuelve contador) \\
\code{def-nppq cond} & Restricción dura: condición nunca verdadera \\
\code{def-tqpq cond} & Restricción dura: condición siempre verdadera \\
\code{def-tdpq cond} & Restricción blanda (como nppq) \\
\code{def-sbqp cond} & Restricción blanda (como tqpq) \\
\code{def-minimizar consulta dominio} & Objetivo: minimizar máximo de consulta \\
\code{def-maximizar consulta dominio} & Objetivo: maximizar mínimo de consulta \\
\code{op-and e1 e2 ...} & Conjunción de expresiones \\
\code{op-or e1 e2 ...} & Disyunción de expresiones \\
\code{op-not e} & Negación \\
\code{op-igual a b} & Igualdad \\
\code{op-distinto a b} & Desigualdad \\
\code{op-mayor a b} & $a > b$ \\
\code{op-menor a b} & $a < b$ \\
\code{(slot entidad)} & Acceso a un atributo de una entidad en el AST \\
\code{generar-python archivo \&key ...} & Genera el archivo Python \\
\bottomrule
\end{tabular}
\end{table}

% ======================================================================
\section{Preguntas frecuentes}
\label{sec:faq}
% ======================================================================

\paragraph{¿Puedo tener más de una entidad combinación?}
No directamente. El algoritmo genético maneja una única lista de
asignaciones. Si el problema requiere múltiples tipos de asignaciones,
considere modelarlo con una sola entidad combinación que tenga más
atributos variables.

\paragraph{¿Qué ocurre si no se encuentran soluciones perfectas?}
El algoritmo genético devuelve la mejor solución encontrada, aunque su
costo sea mayor que cero. Esto puede indicar que el problema es
sobredeterminado (más restricciones duras de las que los datos permiten
satisfacer simultáneamente), que el número de generaciones es insuficiente,
o que los parámetros del algoritmo necesitan ajuste. Los parámetros
\code{POP\_SIZE}, \code{GENERACIONES}, \code{PROB\_CRUCE},
\code{PROB\_MUTACION} y \code{TORNEO\_K} pueden modificarse directamente
en el archivo Python generado.

\paragraph{¿Puedo añadir restricciones manualmente al Python generado?}
Sí. El archivo Python generado es código ordinario que puede editarse.
Sin embargo, se recomienda añadir las restricciones al dominio Lisp y
regenerar, ya que de lo contrario los cambios manuales se perderán en
la próxima generación.

\paragraph{¿Cómo ajusto los pesos de las restricciones?}
A través de los parámetros \code{:peso-dura} y \code{:peso-blanda} de
\code{generar-python}. Si se necesita que ciertas restricciones blandas
tengan mayor prioridad que otras, una alternativa es usar distintos
valores de \code{:peso-blanda} en llamadas separadas (creando dos grupos
de restricciones con pesos distintos) y sumar los costos.

\paragraph{¿El orden de los slots en def-entidad importa?}
Sí, y es crítico. Los argumentos de la macro \code{def-<entidad>} deben
proporcionarse en el mismo orden en que se declararon: primero los
atributos heredados de los mixins (en el orden en que se listaron los
mixins), y luego los atributos propios.

\paragraph{¿Puedo usar el generador sin SBCL?}
El generador está escrito en Common Lisp estándar (ANSI CL). En principio
puede ejecutarse con cualquier implementación de Common Lisp, aunque las
pruebas se han realizado exclusivamente con SBCL.

% ======================================================================
\section*{Conclusión}
\addcontentsline{toc}{section}{Conclusión}
% ======================================================================

El generador de horarios proporciona un flujo de trabajo claro y
reproducible para modelar problemas de asignación:

\begin{enumerate}
  \item \textbf{Definir el dominio} en Lisp: mixins, entidades, combinación,
    consultas, restricciones y datos.
  \item \textbf{Generar el Python} con \code{sbcl --script dominio.lisp}.
  \item \textbf{Ejecutar la metaheurística} con \code{python horario.py}.
  \item \textbf{Analizar el modelo matemático} y el horario propuesto.
\end{enumerate}

La separación entre la descripción del problema (Lisp) y su resolución
(Python generado) permite iterar rápidamente sobre el modelo: modificar
restricciones, añadir datos o cambiar pesos sin necesidad de reescribir
el código de optimización.

\end{document}
